\section{Discussion}
\label{sec:discussion}

\paragraph{Structural hypothesis.}
We hypothesize that the transfer failure stems from structural differences between tool-use and embodied planning. Table~\ref{tab:structural} (Appendix) identifies five dimensions where these domains diverge: error signal (discrete API codes vs.\ diffuse wrong-object errors), observability (explicit stack traces vs.\ silent failures), recovery scope (local retry vs.\ global replanning), feedback timing (immediate vs.\ delayed), and error locality (single-step vs.\ cascading). In tool-use, taxonomy labels encode a direct mapping from error type to recovery action. In embodied planning, errors propagate silently---picking the wrong object produces no error message, and recovery requires replanning from an unknown state. Reflexion \cite{shinn2023reflexion} achieves $\sim$97\% on ALFWorld through inference-time failure reflection, suggesting that inference-time self-correction succeeds where training-time injection fails in embodied planning.

\paragraph{Recovery-pairing confound.}
PALADIN pairs each error with an explicit recovery demonstration via retrieval augmentation, contributing substantially to its +57pp system-level gain. Our TaxEI injects standalone errors without paired recoveries. We cannot distinguish between these two hypotheses with the current design: the negative result may stem from (a)~missing recovery pairing, or (b)~structural differences between tool-use and embodied planning, or both. Recovery-paired injection on ALFWorld is needed to disentangle these factors.

\paragraph{Training lottery at high injection.}
At 50\% injection, individual seeds range from 46.8\% to 70.2\%---the latter exceeding all vanilla seeds. The 70.2\% anomaly occurs in both TaxEI (seed=0) and Random (seed=123), suggesting a generic high-dose phenomenon independent of taxonomy labeling. This cross-condition coincidence may reflect a seed-initialization $\times$ injection interaction warranting further investigation. For practitioners, this implies that single-seed evaluation is unreliable when injection ratios exceed $\sim$25\%: a minimum of three seeds is necessary to distinguish genuine improvements from seed-dependent outliers, and high injection ratios require particular caution.

\paragraph{Detection rate.}
The uniform 0\% detection rate (\S\ref{sec:recovery}) may reflect that SFT does not internalize taxonomy labels when they appear only as prompt prefixes---the training objective (next-token prediction) does not incentivize error classification, so models learn trajectory patterns rather than error type distinctions. This suggests a fundamental limitation of direct SFT mixing for error-awareness: without an explicit classification or contrastive objective, taxonomy labels serve as surface tokens that do not transfer to inference-time error recognition. Methods that use preference optimization \cite{song2024eto} may better incentivize error discrimination.
